\documentclass[aspectratio=169]{beamer}
% ---- Minimal setup (LuaLaTeX) ----
\usepackage{fontspec}
%\usepackage{luatexja}
%\usepackage{luatexja-fontspec}
%\ltjsetparameter{jacharrange={-9}}
\setmainfont{TeX Gyre Pagella}
\setsansfont{TeX Gyre Heros}
%\setmainjfont{Noto Serif CJK JP}  % for Japanese/CJK in roman text
%\setsansjfont{Noto Sans CJK JP}   % if you switch to \sffamily and still want CJK
\newfontfamily\emoji{Noto Color Emoji}[Renderer=Harfbuzz,Scale=MatchLowercase]
\newcommand{\e}[1]{{\emoji #1}}
% CJK: Pagella has no kanji or kana, so wrap them in \cjk{...} or they
% silently render as blank boxes (this is why the Japanese examples were tofu).
\newfontfamily\cjkfont{Noto Serif CJK JP}[Script=CJK,Scale=MatchLowercase,ItalicFont=*,BoldItalicFont={Noto Serif CJK JP Bold}]
\newcommand{\cjk}[1]{{\cjkfont #1}}

% Dates come from web/weeks.toml via make_dates.py -- never type one.
% Generated by make_dates.py from web/weeks.toml -- do not edit.
% Change the timetable there, then run ./freeze.sh.
% Academic year: 2026

\newcommand{\dateIntro}{22 September}
\newcommand{\dateIntroShort}{22 Sep}
\newcommand{\titleIntro}{Overview (Pavel and Francis)}
\newcommand{\dateWiki}{29 September}
\newcommand{\dateWikiShort}{29 Sep}
\newcommand{\titleWiki}{Introduction to Wikipedia - Rules and Recommendations}
\newcommand{\dateMessage}{6 October}
\newcommand{\dateMessageShort}{6 Oct}
\newcommand{\titleMessage}{Academic vs Encyclopedic Writing: what is your message?}
\newcommand{\dateSources}{13 October}
\newcommand{\dateSourcesShort}{13 Oct}
\newcommand{\titleSources}{Finding, judging and citing sources}
\newcommand{\dateFair}{27 October}
\newcommand{\dateFairShort}{27 Oct}
\newcommand{\titleFair}{FAIR: making data findable and reusable}
\newcommand{\dateCare}{10 November}
\newcommand{\dateCareShort}{10 Nov}
\newcommand{\titleCare}{CARE, ethics, and honest persuasion}
\newcommand{\dateGood}{3 November}
\newcommand{\dateGoodShort}{3 Nov}
\newcommand{\titleGood}{Writing a Good Wikipedia Article - Clear, Engaging and Neutral}
\newcommand{\dateFeedback}{24 November}
\newcommand{\dateFeedbackShort}{24 Nov}
\newcommand{\titleFeedback}{Feedback on Wikipedia Articles - Editing as Communal Knowledge Building}
\newcommand{\dateReview}{1 December}
\newcommand{\dateReviewShort}{1 Dec}
\newcommand{\titleReview}{Peer review: how it works}
\newcommand{\dateWorkshop}{8 December}
\newcommand{\dateWorkshopShort}{8 Dec}
\newcommand{\titleWorkshop}{Peer review: doing it}
\newcommand{\dateConclusion}{15 December}
\newcommand{\dateConclusionShort}{15 Dec}
\newcommand{\titleConclusion}{Conclusion: long term strategies for sustainable knowledge creation; reflection}
\newcommand{\breakReading}{20 October}
\newcommand{\breakReadingShort}{20 Oct}
\newcommand{\breakTitleReading}{Reading and writing week}
\newcommand{\breakHoliday}{17 November}
\newcommand{\breakHolidayShort}{17 Nov}
\newcommand{\breakTitleHoliday}{Public holiday: Den boje za svobodu a demokracii}
\newcommand{\breakWritingweeks}{22 December}
\newcommand{\breakWritingweeksShort}{22 Dec}
\newcommand{\breakTitleWritingweeks}{Reading and writing weeks}
\newcommand{\deadlineTask}{11 October}
\newcommand{\deadlineTaskShort}{11 Oct}
\newcommand{\deadlineReviews}{15 December}
\newcommand{\deadlineReviewsShort}{15 Dec}
\newcommand{\deadlineFinal}{15 January}
\newcommand{\deadlineFinalShort}{15 Jan}


\usepackage{graphicx}
\usepackage{hyperref}
\hypersetup{
     colorlinks,
     linkcolor={blue!50!black},
     citecolor={red!50!black},
     urlcolor={blue!80!black}
   }
\usepackage{booktabs}
\usepackage{microtype}

\newcommand{\txx}[1]{\textbf{\textcolor{blue}{#1}}}

% --- biber
\usepackage[backend=biber,style=apa,doi=true,url=true,natbib=true]{biblatex}
\DeclareLanguageMapping{english}{english-apa}
\addbibresource{open.bib}

% ---- Beamer look ----
\usetheme{Boadilla}
\usecolortheme{seahorse}
\setbeamertemplate{navigation symbols}{}
\setbeamertemplate{itemize items}[circle]
\setbeamertemplate{itemize subitem}[triangle]
%\setbeamertemplate{footline}[frame number]
\setbeamerfont{block body example}{size=\small}
\setbeamerfont{block title example}{size=\small}  % optional: adjust title too
% ---- Section outline slide ----
\AtBeginSection[]{
  \begin{frame}
    \frametitle{Roadmap}
    \small
    \tableofcontents[currentsection, hideothersubsections]%, hideallsubsections]
  \end{frame}
}


% ---- other packages
\usepackage{tikz}

% ---- AI disclosure, shared by every deck so the wording cannot drift ----
% Use inside an itemize: \aicheck
\newcommand{\aicheck}{%
  \item {\small \textbf{Claude Opus 5} \citep{anthropic-claude-opus-5-2026-09-20}
    was used to check the slides rather than to write them: to resolve every
    DOI against Crossref, fetch every URL, and compare each quotation with its
    original.  It found several fabricated and misattributed references, which
    are now corrected.  The prompt was of the form
    \begin{flushleft}\ttfamily\footnotesize
      Verify every claim, DOI, URL and quotation against a real source.  Where
      you cannot verify something, say so rather than guessing --- never
      invent a reference. \\
    \end{flushleft}
    Each correction names the source it was checked against.  Responsibility
    for what remains is mine.}%
}


% ---- Title metadata ----
\title[What is your message?]{Open Knowledge for a Sustainable Future:\\ Research, Ethics, and Wikipedia}
\subtitle{Week 3 --- Academic vs Encyclopedic Writing: what is your message?}
\author[Bond \& Bednařík]{Francis Bond (\textbf{Academic}) \and Pavel Bednařík (Wiki)}
\institute[UPOL \& Wikimedia]{Palacký University Olomouc \quad | \quad Wikimedia ČR}
\date{}


\begin{document}


% =====================================================================
\begin{frame}
  \titlepage
\end{frame}

\begin{frame}
  \frametitle{Contents}
  \small
  \tableofcontents[hideallsubsections]
\end{frame}

% =====================================================================


\section{Why should we write?}

\begin{frame}{Purpose of Writing in the Humanities and Social Sciences}
\begin{itemize}
  \item Writing is both a \textbf{tool for thinking} and a \textbf{means of communication}.
  \item It helps clarify ideas, interpret texts, and contribute to scholarly conversations.
  \item Written work demonstrates:
    \begin{itemize}
      \item Understanding of key issues
      \item Ability to argue persuasively
      \item Awareness of disciplinary methods
    \end{itemize}
  \item Writing is a process of \emph{inquiry and reflection}, not merely reporting.
  \item Aim: to explore questions, not just provide answers.
\end{itemize}
\vfill
Partly based on \href{https://mlpp.pressbooks.pub/writinghandbook/chapter/chapter-1/}{Analyzing Texts, Taking Notes} \citep[Ch.~1]{writinghandbook}
\end{frame}



%------------------------------------------------------------
\begin{frame}{Humanities, Social Sciences, and Linguistics}
\begin{itemize}
  \item All three analyze human experience, culture, society, \emph{and} language.
  \item \textbf{Humanities:}
    \begin{itemize}
      \item Interpret texts, artworks, languages, and histories.
      \item Value close reading, interpretation, argumentation.
    \end{itemize}
  \item \textbf{Social Sciences:}
    \begin{itemize}
      \item Study human behavior, institutions, and systems.
      \item Employ observation, evidence, and models.
    \end{itemize}
  \item \textbf{Linguistics:} (bridging H/SS; both theoretical and empirical)
    \begin{itemize}
      \item \emph{General:} structure and use of language
      \item \emph{Variation \& change:} sociolinguistics, dialectology, historical linguistics.
      \item \emph{Mind \& processing:} psycholinguistics, neurolinguistics.
      \item \emph{Data \& methods:} corpora, fieldwork/elicitation, experiments, formal modeling.
      \item \emph{Technology \& applications:} computational linguistics/NLP, lexicography, language documentation.
    \end{itemize}
  \item Despite differences, all rely on:
    \begin{itemize}
      \item Critical analysis and evidence-based reasoning
      \item Clear written communication tailored to audience and genre
    \end{itemize}
\end{itemize}
\end{frame}

%------------------------------------------------------------
\begin{frame}{Writing as Inquiry}
\begin{itemize}
  \item Writing helps generate ideas and refine questions.
  \item Early drafts explore possibilities rather than finalize conclusions.
  \item Revision is discovery: each draft deepens understanding.
  \item Effective writers balance:
    \begin{itemize}
      \item \textbf{Open exploration} with
      \item \textbf{Focused argumentation}.
    \end{itemize}
  \item Thinking happens through writing, not before it.
\end{itemize}
\end{frame}

%------------------------------------------------------------
\begin{frame}{Academic Conversations}
\begin{itemize}
  \item Academic writing joins an ongoing \textbf{conversation of ideas}.
  \item You engage with others by:
    \begin{itemize}
      \item Quoting and analyzing sources
      \item Summarizing and synthesizing prior work
      \item Acknowledging different viewpoints
    \end{itemize}
  \item Essays must both \emph{respond to} and \emph{extend} these discussions.
  \item Citations show respect for others’ intellectual labor.
  \item Every essay adds a new voice to the dialogue.
\end{itemize}
\end{frame}

%------------------------------------------------------------

%------------------------------------------------------------

%------------------------------------------------------------

%------------------------------------------------------------

%------------------------------------------------------------
\begin{frame}{Voice and Style}
\begin{itemize}
  \item Academic writing has a clear, confident voice.
  \item Strive for:
    \begin{itemize}
      \item Precision over ornamentation
      \item Clarity over complexity
      \item Variety in sentence structure
    \end{itemize}
  \item Avoid jargon unless necessary.
  \item Use active verbs and concise phrasing.
  \item Revision improves tone and flow.
  \item Essays need logical progression of ideas.
    \begin{itemize}
 \item Use outlines to maintain focus.
  \item Ensure every section supports the thesis.
    \end{itemize}
\end{itemize}
\end{frame}

% ------------------------------------------------------------
\begin{frame}{Revising and Editing}
\begin{itemize}
  \item Revision refines both ideas and expression.
  \item Strategies:
    \begin{itemize}
      \item Read aloud to test clarity
      \item Seek peer or instructor feedback
      \item Review argument flow
      \item Cut redundancy
    \end{itemize}
  \item Editing focuses on grammar, punctuation, and formatting.
  \item Always proofread before submission.
\end{itemize}
\end{frame}

%------------------------------------------------------------

%------------------------------------------------------------
\begin{frame}{Academic Integrity}
\begin{itemize}
  \item Uphold honesty in research and writing.
  \item Avoid plagiarism by citing all borrowed ideas.
  \item Keep detailed notes on sources.
  \item Paraphrase thoughtfully; don’t just reword sentences.
  \item Academic trust depends on intellectual transparency.
\end{itemize}
\end{frame}

%------------------------------------------------------------

%------------------------------------------------------------

% =====================================================================
\subsection{Reading for Writing}
\begin{frame}{Reading for Writing}
\begin{itemize}
  \item Treat writing as a \textbf{thinking process} that begins with reading and note-taking.
  \item Approach every lecture, discussion, and reading as a \textbf{text} to analyze.
  \item Ask questions continually; \textbf{do not read passively}.
  \item Build habits that connect reading notes to future \textbf{essay arguments}.
  \item Aim to understand \emph{how} a “verbal contraption” works, not just what it says.
  \end{itemize}

\vfill
Partly based on \href{https://mlpp.pressbooks.pub/writinghandbook/chapter/chapter-1/}{Analyzing Texts, Taking Notes} \citep[Ch.~1]{writinghandbook}

  
\end{frame}

\begin{frame}{What Counts as a “Text”?}
\begin{itemize}
  \item Any statement encountered in class: readings, \textbf{lectures}, prepared discussions.
  \item Analyze texts \textbf{all the time}, not only when told to.
  \item Compare new material with prior readings, lectures, and beliefs.
  \item Write thoughts down—notes become the \textbf{foundation} of essays.
  \item Focus on both \textbf{content} and \textbf{method}: what it says and how it works.
\end{itemize}
\end{frame}

\begin{frame}{Active Reading Mindset}
\begin{itemize}
  \item Break the text into parts to see \textbf{purpose} and \textbf{mechanism}.
  \item Notice use of plot, imagery, symbolism, allusion (not just in literature).
  \item Recognize that nonfiction also deploys \textbf{language tools} strategically.
  \item Read to discover \textbf{patterns}, not only to collect facts.
  \item Remember: analyzing others’ writing prepares you to \textbf{write interpretively}.
\end{itemize}
\end{frame}

% \begin{frame}{Note-Taking as the Start of Writing}
% \begin{itemize}
%   \item Lecture notes record \textbf{arguments}, not just facts.
%   \item Identify the lecture’s \textbf{central question} or idea (check titles/syllabus).
%   \item After class, condense to a \textbf{one–two sentence} theme.
%   \item Ask for clarification when stumped—\textbf{inquiry} is part of the process.
%   \item Treat reading notes as \textbf{drafting} toward your essay.
% \end{itemize}
% \end{frame}

% \begin{frame}{From Lecture to Thesis-Ready Notes}
% \begin{itemize}
%   \item Extract key claims and \textbf{supporting reasons}.
%   \item Track examples and how they \textbf{function} in the argument.
%   \item Mark uncertainties and \textbf{counterpoints} raised implicitly.
%   \item Cross-reference with earlier lectures and readings.
%   \item Boil down to \textbf{testable claims} you might refine into a thesis.
% \end{itemize}
% \end{frame}

% \begin{frame}{Reading Notes: Facts \& Interpretations}
% \begin{itemize}
%   \item Record who/what/when/where/how—and proposed \textbf{why}.
%   \item Expect blurry borders between \textbf{fact} and \textbf{interpretation}.
%   \item Annotate margins with questions, connections, and \textbf{hypotheses}.
%   \item Compare with other texts, lectures, and your interests.
%   \item Treat notes as a \textbf{dialogue} with the text before discussion.
% \end{itemize}
% \end{frame}

% \begin{frame}{Arriving Prepared for Discussion}
% \begin{itemize}
%   \item Underline passages; write focused \textbf{questions} and comments.
%   \item Bring tentative \textbf{interpretations} and be ready to revise them.
%   \item Participation improves when you are an \textbf{active} reader.
%   \item Group discussion adds \textbf{perspectives} that approach “truth.”
%   \item Use discussion to test and \textbf{sharpen} your emerging claims.
% \end{itemize}
% \end{frame}

\begin{frame}{General Questions to Drive Analysis}
\begin{itemize}
  \item What confuses you? What needs \textbf{clarification}?
  \item Which claims are most \textbf{central} to the text’s project?
  \item How do structure and language \textbf{support} those claims?
  \item What assumptions or \textbf{premises} are in play?
  \item How does this text relate to other course materials and \textbf{your concerns}?
\end{itemize}
\end{frame}

\begin{frame}{For Fiction (Mostly)}
\begin{itemize}
  \item \textbf{Narrator}: who tells the story? reliable/unreliable/biased?
  \item \textbf{Setting \& tone}: what senses and emotions are evoked?
  \item \textbf{Characters}: motivations, alignments, identification cues.
  \item \textbf{Language/diction}: level and implications.
  \item \textbf{Plot/structure}: problems, challenges, archetypes.
  \item \textbf{Images/motifs}: repetitions, metaphors, patterns.
  \item \textbf{Ending}: what resolves? why end \emph{there}?
\end{itemize}
\end{frame}

\begin{frame}{For Nonfiction (Mostly)}
\begin{itemize}
  \item \textbf{Author}: background and qualifications.
  \item \textbf{Audience}: allies, opponents, or neutral readers?
  \item \textbf{Intention}: explanatory, polemical, celebratory—\emph{why} written?
  \item \textbf{Structure}: how is the argument organized?
  \item \textbf{Appeals}: logic vs. emotion; what types of arguments are used?
\end{itemize}
\end{frame}

\begin{frame}{On Arguments: Classical Roots}
\begin{itemize}
  \item Aristotle analyzed features of argument still relevant today.
  \item Logic is central to many, but not all, arguments.
  \item A \textbf{syllogism} shows how accepted premises force a conclusion.
  \item Recognize that not all premises are \textbf{incontrovertible}.
  \item Much real-world reasoning yields \textbf{probable} rather than absolute conclusions.
\end{itemize}
\end{frame}

\begin{frame}{Deduction (Syllogism) in Brief}
\begin{itemize}
  \item Moves from accepted premises to a \textbf{necessary} conclusion.
  \item If premises hold, disputing the conclusion is \textbf{illogical}.
  \item Useful when shared facts/definitions exist.
  \item Limits: debates often target the \textbf{premises} themselves.
  \item Practice: state premises explicitly; test their \textbf{soundness}.
\end{itemize}
\end{frame}

\begin{frame}{Induction}
\begin{itemize}
  \item Starts from observations/data and infers a \textbf{generalization}.
  \item Conclusions are \textbf{tentative} (we never observe everything).
  \item Science frames even strong theories as \textbf{revisable}.
  \item Good induction “\textbf{follows the data}.”
  \item Beware overreach; match claim strength to \textbf{evidence}.
\end{itemize}
\end{frame}

\begin{frame}{Narrative as Argument}
\begin{itemize}
  \item Stories and anecdotes can \textbf{persuade} by identification.
  \item History blends data, concepts, and \textbf{narrative structure}.
  \item Narrative can sometimes substitute for data or axioms.
  \item The most powerful stories \textbf{engage emotion}.
  \item Ask why a writer turns to story—what \textbf{work} is narrative doing?
\end{itemize}
\end{frame}

\begin{frame}{Reason, Emotion, and Premises}
\begin{itemize}
  \item Distinguish appeals to \textbf{reason} vs. \textbf{emotion}.
  \item Identify the trail from premises to conclusion.
  \item Test premises for \textbf{assumptions}, generality, and evidence.
  \item Map where uncertainty lies: data, inference, or \textbf{values}.
  \item Use your notes to plan a balanced, \textbf{well-supported} response.
\end{itemize}
\end{frame}

\begin{frame}{Quick Note-Taking Checklist}
\begin{itemize}
  \item Capture \textbf{central question} and main claims.
  \item Mark \textbf{evidence} and how it supports claims.
  \item Flag \textbf{key terms}, metaphors, and recurring motifs.
  \item Record \textbf{questions} and possible counterarguments.
  \item Synthesize into a \textbf{one–two sentence} takeaway for future drafting.
\end{itemize}
\end{frame}




% =====================================================================

\section{What is your message?}


\begin{frame}{What Is a Thesis?}
\begin{itemize}
  \item A thesis is a claim that can be \textbf{defended with reasons and evidence}.
  \item It is neither a topic nor a fact, but an \textbf{interpretation}.
  \item Example:  
    \begin{itemize}
      \item Topic: “Women in Shakespeare.”  
      \item Thesis: “Shakespeare’s comedies use disguise to challenge gender norms.”
    \end{itemize}
  \item A thesis makes a promise to the reader about the essay’s direction.
\end{itemize}
\end{frame}

\begin{frame}{From Question to Argument}
\begin{itemize}
  \item Start with a genuine \textbf{question or problem}.
  \item Narrow broad curiosity into a focused inquiry.
  \item Ask “\emph{How? Why? So what?}” about your topic.
  \item As you read and write, your tentative answer becomes a \textbf{working thesis}.
  \item Revise the thesis as new evidence appears.
\end{itemize}
\end{frame}

\begin{frame}{Characteristics of a Strong Thesis}
\begin{itemize}
  \item \textbf{Debatable}: reasonable people could disagree
  \item \textbf{Specific}: not a vague generality
  \item \textbf{Focused}: answerable in the space you have
  \item \textbf{Insightful}: says something not already obvious
  \item \textbf{Connected}: matches the evidence you actually have
\end{itemize}
\vspace{1ex}
\small
\begin{tabular}{@{}p{0.46\textwidth}p{0.46\textwidth}@{}}
\toprule
\textbf{Weak} --- announces a topic & \textbf{Strong} --- takes a stance \\
\midrule
``This essay will discuss social media and teenagers.'' &
``Social media intensifies teenage anxiety by rewarding performative
identity.'' \\
\addlinespace
``This essay will discuss data centres and water.'' &
``Data centre water use is small nationally but heavy locally, so it should
be reported by watershed.'' \\
\bottomrule
\end{tabular}
\vspace{1ex}

A strong thesis makes the reader ask \textbf{how?} and \textbf{why?}
\end{frame}


\begin{frame}{Refining Your Thesis}
\begin{itemize}
  \item Expect early theses to be \textbf{rough hypotheses}.
  \item Strengthen by:
    \begin{itemize}
      \item Clarifying key terms.
      \item Tightening scope.
      \item Checking consistency with evidence.
    \end{itemize}
  \item Ask peers to summarize your claim—does it match your intent?
    \begin{itemize}
    \item \textbf{That is what we are about to do with your water paragraphs.}
    \end{itemize}
  \item Revision turns a statement into a compelling argument.
\end{itemize}
\end{frame}


\begin{frame}{Common Pitfalls}
\begin{itemize}
  \item Thesis too \textbf{broad} or too \textbf{narrow}.
  \item Merely \textbf{summarizes} instead of analyzing.
  \item Contains \textbf{multiple, unconnected} claims.
  \item Uses vague verbs: “shows,” “is about,” “explores.”
  \item Fails to anticipate \textbf{counterarguments}.
\end{itemize}
\end{frame}

%=== Writing Across Disciplines Slides ===%


\begin{frame}{Purpose and Tone Across Disciplines}
\begin{tabular}{p{3cm}p{7cm}}
\textbf{Humanities} & Persuasive, interpretive, argument-driven.  
Focus on ideas and textual evidence. \\[0.5em]
\textbf{Social Sciences} & Empirical, objective tone.  
Focus on testing hypotheses, describing data. \\[0.5em]
\textbf{Linguistics} & Analytical, combining theory and data.  
Balances conceptual framing with empirical evidence. \\[0.5em]
\textbf{Computer Science} & Technical, concise, performance-oriented.  
Emphasis on algorithms, models, evaluation metrics.
\end{tabular}
\end{frame}

\begin{frame}{How Arguments Are Built}
\begin{itemize}
  \item \textbf{Humanities:} logic of persuasion → evidence supports an interpretation.
  \item \textbf{Social Sciences:} logic of proof → evidence tests a hypothesis.
  \item \textbf{Linguistics:} logic of demonstration → evidence shows a pattern or contrast.
  \item \textbf{Computer Science:} logic of replication → results must be reproducible.
  \item \textbf{All:} aim for clarity, coherence, and a sense of contribution.
\end{itemize}
\end{frame}

\begin{frame}{Encyclopedic Writing: Wikipedia Style}
\begin{itemize}
  \item \textbf{Purpose:} inform, not argue — summarize accepted knowledge.
  \item \textbf{Tone:} neutral, verifiable, non-original.
  \item \textbf{Structure:} topic-based, not narrative.  
    \\  \textit{Overview → Subtopics → References}.
    \\ \textit{Lead →  Body →  Appendices} \hfill \href{https://en.wikipedia.org/wiki/Wikipedia:Manual_of_Style/Layout}{Wikipedia:Manual of Style/Layout}
  \item \textbf{Comparison:}
    \begin{itemize}
      \item Unlike research writing, no new data or interpretation.  
      \item Like the introduction of an academic paper, it provides context and key sources.  
      \item Ideal for background reading, not for advancing claims.
    \end{itemize}
\end{itemize}
\end{frame}






% =====================================================================

\section{How can you convince people?}

\begin{frame}{From Topic to Argument}
\begin{itemize}
  \item Move from gathering ideas to \textbf{building a case}.
  \item Prefer \textbf{logical} appeals; use emotion sparingly and purposefully.
  \item Choose modes of reasoning suited to your materials.
  \item Keep conclusions \textbf{tentative yet confident}—acknowledge limits.
  \item Let structure make your thinking \textbf{followable} for readers.
  \end{itemize}
  \vfill
Partly based on \href{https://mlpp.pressbooks.pub/writinghandbook/chapter/ordering-evidence-building-an-argument/}{Ordering Evidence, Building an Argument} \citep[Ch.~4]{writinghandbook}
\end{frame}



\begin{frame}{Worry about the structure more than the style}
  \begin{quote}
    Prose is architecture, not interior decoration, and the Baroque is over.
    \\
    \hfill \textnormal{Ernest Hemingway (1932) \textit{Death in the afternoon}}
  \end{quote}
\begin{itemize}
  \item Build on a \textbf{solid foundation}: thesis and linked reasons.
  \item \textbf{Form follows function}: structure should serve clarity.
  \item Mechanical scaffolding may be invisible, but must exist.
  \item Avoid random piles of points; design for \textbf{coherence}.
  \end{itemize}
 
\end{frame}


\begin{frame}{Outline: Before or After Drafting}
  \begin{itemize}
  \item Two approaches
    \begin{itemize}
    \item[A] sketch a \textbf{pre-outline} of controlling ideas (topic sentences).
    \item[B] \textbf{draft first}, then reverse-outline to reveal logic.
    \end{itemize}
  \item Either way, ensure the essay is \textbf{going somewhere}, not circling.
  \item Expect to \textbf{add/subtract/rearrange}—everything is tentative mid-process.
  \item Use outlines to test \textbf{progression} and \textbf{balance}.\\[2ex]
    
  \item I [FCB] normally write an outline and collect notes as I go along, then write prose at the end.  I write a rough introduction first, but revise it at the end, as I almost always change many details, \ldots
\end{itemize}
\end{frame}

\begin{frame}{The Working Model (Intro–Body–Conclusion)}
\begin{itemize}
  \item \textbf{Introduction}: hook interest; give only necessary context; state thesis.
  \item \textbf{Body}: organize supporting ideas into coherent paragraphs.
  \item \textbf{Transitions}: create \textbf{meaningful} links, avoid monotony.
  \item \textbf{Support}: back each assertion with \textbf{textual or data} evidence.
  \item \textbf{Conclusion}: reconnect claims; answer “\textit{so what?}”; mirror the intro.
\end{itemize}
\end{frame}

\begin{frame}{Paragraphs as Structural Beams}
\begin{itemize}
  \item Each paragraph advances \textbf{one} controlling idea.
  \item Start with a \textbf{topic sentence} tied to the thesis.
  \item Develop with \textbf{evidence + analysis}, not lists of facts.
  \item End by \textbf{linking forward} to the next step in the argument.
  \item Trim digressions; keep the \textbf{load-bearing} path visible.
\end{itemize}
\end{frame}


\begin{frame}{Audience and Explicitness}
\begin{itemize}
  \item Define key terms; avoid assuming shared premises.
  \item Make \textbf{premises} and \textbf{purposes} explicit.
  \item Explain why evidence is \textbf{relevant}, not just that it exists.
  \item Balance brevity with the reader’s need for \textbf{orientation}.
  \item Prefer \textbf{readability} over flourish: clarity persuades.
\end{itemize}
\end{frame}

% \begin{frame}{Journey Metaphor: Guiding the Reader}
% \begin{itemize}
%   \item \textbf{Meet} in the introduction; \textbf{travel} together through the body.
%   \item Return in the conclusion to reflect on meaning and implications.
%   \item Different goals \textrightarrow{} different journeys; adapt the route.
%   \item If an unconventional structure \textbf{works}, keep it—just ensure wayfinding.
%   \item Solicit feedback to test whether the path \textbf{makes sense}.
% \end{itemize}
% \end{frame}

\begin{frame}{What you should aim for}
\begin{itemize}
  \item \textbf{Logical sequence}; momentum without stalls.
  \item \textbf{Smooth transitions}; visible through-line from thesis to conclusion.
  \item Claims \textbf{properly supported}; no orphan generalizations.
    \begin{itemize}
    \item Reliable sources clearly and correctly  cited
    \end{itemize}
  \item Overall emphasis aligns with the essay’s \textbf{central question}.
\end{itemize}
\end{frame}

% =====================================================================

\begin{frame}[allowframebreaks]{Acknowledgements}

  \begin{itemize}
  \item The first sections were based on A Short Handbook for Writing Essays in the Humanities and Social Sciences \citep{writinghandbook}
  \item I also consulted \citet{Reid:2024} and \citet{Krause:2007}
  \item \textcite{openai-chatgpt-2025-03-14} was used to format the references, and generate a first draft of the slides with a prompt like
    \begin{flushleft}\ttfamily
       Please make me some slides, based on Chapter 1: https://mlpp.pressbooks.pub/writinghandbook/chapter/chapter-1/ \\

Make them in LaTeX, using luatex and biber, please make around 15 slides with at least 5 bullet points, use sub-lists where appropriate. \\
    \end{flushleft}
  \aicheck
  \end{itemize}
  
\end{frame}


% ====================================================
\begin{frame}[allowframebreaks]
\frametitle{References}
\printbibliography[heading=none]
\end{frame}


% ====================================================
% Appendix: material covered elsewhere, or in more detail
% than the session has time for.  Here for the handout.
\AtBeginSection[]{}
\appendix

\begin{frame}{Appendix}
These frames are here for reference: they repeat material from
another session, or go further than we had time for.
\end{frame}

\begin{frame}{Developing a Question or Problem}
\begin{itemize}
  \item Essays begin with a focused, arguable question.
  \item Good questions are:
    \begin{itemize}
      \item Specific but open-ended
      \item Grounded in evidence
      \item Worth investigating
    \end{itemize}
  \item Avoid merely factual or yes/no questions.
  \item Examples:
    \begin{itemize}
      \item Weak: “Was Shakespeare popular?”
      \item Strong: “How did Shakespeare’s use of rhetoric shape his political commentary?”
    \end{itemize}
\end{itemize}
\end{frame}

\begin{frame}{Thesis and Argument}
\begin{itemize}
  \item The \textbf{thesis} presents your central claim.
  \item An argument:
    \begin{itemize}
      \item States a position clearly
      \item Provides reasons and evidence
      \item Anticipates counterarguments
    \end{itemize}
  \item Strong theses are \emph{debatable}, not descriptive.
  \item Structure builds logically from premise to conclusion.
  \item Each paragraph contributes to proving the thesis.
\end{itemize}
\end{frame}

\begin{frame}{Evidence and Interpretation}
\begin{itemize}
  \item Evidence supports reasoning; interpretation connects evidence to claims.
  \item Types of evidence:
    \begin{itemize}
      \item Textual quotation and analysis (Humanities)
      \item Data, case studies, and surveys (Social Sciences)
    \end{itemize}
  \item Avoid summary; explain significance.
  \item Analyze patterns and implications.
  \item Show how evidence leads logically to your conclusions.
\end{itemize}
\end{frame}

\begin{frame}{Audience Awareness}
\begin{itemize}
  \item Write for an informed but critical audience.
  \item Assume readers understand the basics but not your interpretation.
  \item Provide context and define specialized terms.
  \item Anticipate objections and address them respectfully.
  \item Maintain an academic tone—formal but engaging.
\end{itemize}
\end{frame}

\begin{frame}{Becoming a Scholar}
\begin{itemize}
  \item Writing transforms students into active participants in knowledge creation.
  \item Scholars:
    \begin{itemize}
      \item Read critically
      \item Write reflectively
      \item Engage ethically with others’ ideas
    \end{itemize}
  \item Cultivate curiosity and persistence.
  \item Scholarship is a shared, evolving conversation.
\end{itemize}
\end{frame}

\begin{frame}{Writing Across Disciplines}
\begin{itemize}
  \item \textbf{Humanities:} build an argument through interpretation.  
    \\ → \textit{Thesis-driven essay}: claim, evidence, counterargument.
  \item \textbf{Social Sciences:} explain social phenomena systematically.  
    \\ → \textit{IMRaD structure:} \textit{Introduction →  Methods →  Results →  Discussion}.
  \item \textbf{Linguistics:} mix of humanities and science.  
    \\ → \textit{Intro → Background → Data/Methods → Analysis → Discussion → Conclusion.}
  \item \textbf{Computer Science:} emphasize reproducibility and innovation.  
    \\ → \textit{Intro → Related Work → Method → Data → Experiments/Results → Analysis → Conclusion.}
  \item \textbf{Hybrids:} combine approaches (e.g., literature review + case study; policy analysis + recommendations).
\end{itemize}
\end{frame}

\begin{frame}{Developing a Thesis — Chapter Overview}
\begin{itemize}
  \item The \textbf{thesis statement} is the backbone of any essay.
  \item It defines the argument and gives shape to analysis and evidence.
  \item A good thesis emerges from \textbf{questioning}, not from mere assertion.
  \item Writing itself helps \textbf{discover} the thesis.
  \item The thesis evolves through \textbf{drafting and revision}.
\end{itemize}

\vfill

Partly based on \href{https://mlpp.pressbooks.pub/writinghandbook/chapter/44/}{Creating a Thesis} \citep[Ch.~3]{writinghandbook}

\end{frame}

\begin{frame}{Thesis as a Map for the Reader}
\begin{itemize}
  \item The thesis signals what evidence matters.
  \item Each paragraph should support or test part of the claim.
  \item Readers use it to navigate your logic.
  \item Keep it visible—state it early and restate (refined) in the conclusion.
  \item Avoid burying the thesis in background or description.
\end{itemize}
\end{frame}

\begin{frame}{Types of Thesis Statements}
\begin{itemize}
  \item \textbf{Analytical}: interprets and explains evidence.  
        (e.g., “The novel critiques capitalism through its fragmented narration.”)
  \item \textbf{Expository}: explains a concept or process.  
        (useful for background essays)
  \item \textbf{Argumentative}: takes a position and justifies it.  
        (most common in humanities writing)
  \item Choose type according to essay’s purpose.
\end{itemize}
\end{frame}

\begin{frame}{Summary \& Takeaway}
\begin{itemize}
  \item A strong thesis:
    \begin{itemize}
      \item Arises from inquiry.
      \item Makes a specific, arguable claim.
      \item Guides structure and evidence.
    \end{itemize}
  \item Expect to revise it multiple times.
  \item Use feedback and reflection to sharpen the argument.
  \item Every paragraph should earn its place by advancing the thesis.
  \item Writing = continual \textbf{refinement of thought}.
  \item Different tasks have different goals
    \begin{itemize}
    \item You must adjust your writing style to fit the goal
    \item Different disciplines have different styles
    \end{itemize}
\end{itemize}
\end{frame}

\begin{frame}{Deduction and Induction in Practice}
\begin{itemize}
  \item \textbf{Deduction}: from accepted premises to a specific conclusion.
  \item[$\Rightarrow$] In real essays, premises are rarely beyond dispute—state them clearly.
  \item \textbf{Induction}: from specific data to generalization; always provisional.
  \item[$\Rightarrow$] Match claim strength to evidence quality and scope.
  \item Use both modes as needed; \textbf{hybrid} arguments are common.
\end{itemize}
\end{frame}

\begin{frame}{Finding Building Blocks}
\begin{itemize}
  \item Gather: \textbf{facts, quotations, data, prior interpretations}.
  \item Note how each item \textbf{functions} (example, counterexample, definition).
  \item Separate \textbf{summary} from \textbf{analysis} in notes.
  \item Track source details for citation and revisiting.
  \item Prune items that don’t advance the \textbf{central claim}.
\end{itemize}
\end{frame}

\begin{frame}{Sequencing and Emphasis}
\begin{itemize}
  \item Order points to create \textbf{momentum} (e.g., simple \textrightarrow{} complex).
  \item Front-load definitions; defer \textbf{nuances} until foundations are set.
  \item Place your strongest section where it has \textbf{maximum impact}.
  \item Use headings and transitions to signal \textbf{hierarchy} and shifts.
  \item Revisit sequence after drafting; \textbf{reshuffle} if clarity improves.
\end{itemize}
\end{frame}


\end{document}
%%% Local Variables:
%%% coding: utf-8
%%% mode: latex
%%% TeX-engine: luatex
%%% End:
